\chapter{Verificación con Coq}
\label{cap:coq}

Mientras que TLA\textsuperscript{+} explora modelos finitos por
enumeración, Coq es un \emph{asistente de pruebas} que verifica
teoremas universalmente cuantificados de manera inductiva. La
verificación con Coq en sotFS apunta a probar que cada regla DPO del
cap.~\ref{cap:dpo} preserva la well-formedness del TypeGraph, sin
acotar el tamaño del grafo.

Este capítulo explica la estructura de los archivos Coq, los teoremas
centrales, y --- crítico para la honestidad del proyecto --- los tres
lemmas \texttt{Admitted} que quedan pendientes en \citaCoq{DpoRmdir.v}.

\section{El modelo Records: \texttt{SotfsGraph.v}}

\citaCoq{SotfsGraph.v} (585 LOC) define el universo formal sobre el
que se prueban las reglas. Las estructuras centrales:

\begin{lstlisting}[language=coq, caption={Records principales en
\texttt{SotfsGraph.v}.}, label={lst:coq-records}]
Record InodeRec := {
  ir_id : InodeId;
  ir_vtype : VnodeType;
  ir_link_count : nat
}.

Record DirRec := {
  dr_id : DirId;
  dr_inode : InodeId
}.

Record ContainsEdge := {
  ce_dir : DirId;
  ce_ino : InodeId;
  ce_name : Name
}.

Record Graph := {
  g_inodes : list InodeRec;
  g_dirs : list DirRec;
  g_edges : list ContainsEdge
}.
\end{lstlisting}

\paragraph{Por qué Records (no inductivos).} Records permiten una
representación operacional muy cercana al runtime Rust: el grafo es
una tripleta de listas. Las pruebas resultan más concretas (less
case-splitting sobre constructores) y la traducción mental con el
código es directa.

\subsection*{Well-formedness}

\begin{lstlisting}[language=coq]
Definition WellFormed (g : Graph) : Prop :=
  TypeInvariant g
  /\ LinkCountConsistent g
  /\ UniqueNamesPerDir g
  /\ NoDanglingEdges g
  /\ NoDirCycles g.
\end{lstlisting}

Cinco cláusulas; las dos invariantes adicionales del runtime
(\texttt{DirHasSelfRef}, \texttt{BlockRefcountConsistent}) no están
formalizadas en el modelo Coq de v0.2.5 porque viven en el subsistema
de bloques que no aparece en estos archivos. La extensión es trabajo
abierto.

\section{Los teoremas de preservación}

Para cada operación POSIX (excepto \opwrite, todavía no formalizado
en Coq), el archivo \citaCoq{DpoX.v} prueba seis teoremas: uno por
cada cláusula del \texttt{WellFormed} y uno que los combina.

\begin{table}[h]
\centering
\renewcommand{\arraystretch}{1.15}
\begin{tabularx}{\linewidth}{l r X}
  \toprule
  \textbf{Archivo} & \textbf{LOC} & \textbf{Teorema final} \\
  \midrule
  \citaCoq{DpoCreate.v} & 348 & \texttt{create\_preserves\_WellFormed} (línea 335) \\
  \citaCoq{DpoMkdir.v}  & 586 & \texttt{mkdir\_preserves\_WellFormed}  (línea 572) \\
  \citaCoq{DpoLink.v}   & 382 & \texttt{link\_preserves\_WellFormed}   (línea 369) \\
  \citaCoq{DpoUnlink.v} & 314 & \texttt{unlink\_keep\_preserves\_WellFormed} (línea 300) \\
  \citaCoq{DpoRmdir.v}  & 513 & \texttt{rmdir\_preserves\_WellFormed}  (línea 492) \\
  \citaCoq{DpoRename.v} & 317 & \texttt{rename\_preserves\_WellFormed} (línea 304) \\
  \bottomrule
\end{tabularx}
\caption{Los seis archivos de prueba DPO. Para cada uno, el teorema
final es de la forma ``$\forall g, \text{precond} \Rightarrow
\text{WellFormed}(g) \Rightarrow \text{WellFormed}(\text{op}\,g)$''.}
\label{tab:coq-files}
\end{table}

\subsection*{Estructura de una prueba típica}

Tomemos \texttt{create\_preserves\_WellFormed} en \citaCoq{DpoCreate.v}
línea 335 como ejemplo. La forma:

\begin{lstlisting}[language=coq]
Theorem create_preserves_WellFormed :
  forall g d n uid gid,
    WellFormed g ->
    name_unique_in g d n ->
    contains_dir g d ->
    WellFormed (create g d n uid gid).
Proof.
  intros g d n uid gid Hwf Huniq Hd.
  destruct Hwf as [Hti [Hlcc [Hun [Hnd Hndc]]]].
  unfold WellFormed; repeat split.
  - apply create_preserves_TypeInvariant; assumption.
  - apply create_preserves_LinkCountConsistent;
    assumption.
  - apply create_preserves_UniqueNamesPerDir;
    assumption.
  - apply create_preserves_NoDanglingEdges; assumption.
  - apply create_preserves_NoDirCycles; assumption.
Qed.
\end{lstlisting}

La descomposición es modular: el teorema final invoca cinco lemmas
auxiliares, uno por cláusula del \texttt{WellFormed}, cada uno probado
arriba en el mismo archivo. Esto hace el código de prueba navegable y
re-usable entre operaciones (\texttt{TypeInvariant} es similar entre
\opcreate{} y \oplink, por ejemplo).

\section{Deuda formal: los tres \texttt{Admitted}}

\begin{trampa}[Tres lemmas \texttt{Admitted} en \texttt{DpoRmdir.v}]
\label{tr:coq-admitted}
En la versión \texttt{v0.2.5}, el archivo \citaCoq{DpoRmdir.v} contiene
tres ocurrencias de \texttt{Admitted} (líneas 349, 405, 505). Todas
están bloqueadas por la misma deficiencia del modelo: la propiedad de
que \emph{los directorios no pueden ser hard-linked} (\texttt{GC-LINK-2}
en la nomenclatura del proyecto) no está enunciada como un lemma
auxiliar de \texttt{WellFormed}, lo que impide cerrar la prueba en los
sub-casos donde habría que saber que el inode removido no es target de
ninguna otra arista \econtains{} sobreviviente.

\textbf{Cómo se ve.} Los tres \texttt{Admitted} aparecen al final de
los sub-casos donde el contexto requiere algo como
\texttt{forall ce \textbackslash in g.edges, ce.ino <> removed\_inode}.
Sin un invariante explícito \texttt{NoHardLinkToDir}, esa hipótesis no
está en scope; el comentario adyacente lo dice:

\begin{lstlisting}[language=coq]
(* NOTE: The proof is complete except for one
   sub-case: showing that no surviving edge targets
   the removed inode ti. This requires the
   additional fact that directories cannot be
   hard-linked (GC-LINK-2), which is not yet part
   of WellFormed. We state the theorem with
   Admitted for transparency. *)
\end{lstlisting}

\textbf{Status declarado.} Los teoremas
\texttt{rmdir\_preserves\_WellFormed} y
\texttt{rmdir\_preserves\_NoDanglingEdges} \emph{no están probados}
en la versión actual. Las pruebas \emph{están casi completas}; el
gap es estructural.

\textbf{Roadmap.} Tres pasos:
\begin{enumerate}
  \item Extender \texttt{WellFormed} con una sexta cláusula
    \texttt{NoHardLinkToDir : Prop} que diga ``ningún directorio es
    target de más de una arista \econtains{} viva (excluyendo
    \texttt{..})''.
  \item Re-probar \texttt{create\_preserves\_WellFormed},
    \texttt{mkdir\_preserves\_WellFormed},
    \texttt{link\_preserves\_WellFormed} y los demás incluyendo el
    nuevo invariante. Es trabajo mecánico: las reglas ya respetan
    la propiedad, sólo falta probarlo formalmente.
  \item Cerrar los tres \texttt{Admitted} usando
    \texttt{NoHardLinkToDir} como hipótesis auxiliar.
\end{enumerate}

\textbf{Lección.} La verificación mecánica obliga a hacer explícitos
invariantes que el código respeta ``por construcción'' pero que el
modelo formal no enuncia. Los \texttt{Admitted} son la deuda visible;
cerrar la deuda significa hacer crecer el modelo, no el código.
\end{trampa}

\section{Cómo correr las pruebas Coq}

\begin{lstlisting}[language=shellcmd]
cd formal/coq
coq_makefile -f _CoqProject -o Makefile
make
\end{lstlisting}

La compilación verifica todos los teoremas; un \texttt{Admitted}
\emph{no} interrumpe la compilación (Coq sí lo loguea como warning).
Para ver los Admitteds y otros warnings:

\begin{lstlisting}[language=shellcmd]
make 2>&1 | grep -E "Admitted|Warning"
\end{lstlisting}

Tiempo total de compilación en una CPU moderna: $\sim$\,2 minutos.

\section{Lo que Coq prueba (y lo que no)}

\paragraph{Lo que prueba.} Para cada regla DPO (salvo
\texttt{rmdir\_preserves\_WellFormed}, \texttt{rmdir\_preserves\_NoDanglingEdges},
y \texttt{rmdir\_preserves\_TypeInvariant}), que la aplicación de la
regla preserva las cinco cláusulas del \texttt{WellFormed} para
\emph{todo} grafo y \emph{todos} los parámetros admisibles.

\paragraph{Lo que no prueba.}
\begin{itemize}
  \item \texttt{rmdir} con sus tres \texttt{Admitted}.
  \item Los invariantes del runtime que no están en el modelo Coq
    (\texttt{DirHasSelfRef}, \texttt{BlockRefcountConsistent},
    \texttt{NoHardLinkToDir}).
  \item La correspondencia entre las funciones \texttt{create},
    \texttt{mkdir}, \dots de Coq y las funciones Rust de
    \citaCrate{sotfs-ops}. Esa correspondencia es por inspección
    humana; no hay extracción Coq~$\to$~Rust.
  \item Las propiedades concurrentes (RCU, GTXN). Las pruebas Coq
    asumen serial; concurrencia se cubre con TLA\textsuperscript{+} y
    tests \texttt{loom}.
\end{itemize}

\section*{Resumen del capítulo}

\begin{itemize}
  \item Coq verifica que cada regla DPO preserva
    \texttt{WellFormed} para grafos arbitrarios; complementa la
    enumeración finita de TLC con prueba inductiva.
  \item \citaCoq{SotfsGraph.v} define los Records \texttt{InodeRec},
    \texttt{DirRec}, \texttt{ContainsEdge}, \texttt{Graph} y la
    proposición \texttt{WellFormed} con cinco cláusulas.
  \item Seis archivos \texttt{Dpo*.v} prueban
    \texttt{X\_preserves\_WellFormed} para create, mkdir, link, unlink,
    rmdir, rename.
  \item \textbf{Deuda formal explícita}: tres \texttt{Admitted} en
    \texttt{DpoRmdir.v} bloqueados por el invariante auxiliar
    \texttt{NoHardLinkToDir} faltante.
  \item Limitaciones: el modelo Coq omite varios invariantes del
    runtime y no hay extracción Coq $\to$ Rust.
\end{itemize}

\section*{Ejercicios}

\begin{ejercicio}[\dificultad{$\star$}]
\label{ej:coq-1}
¿Cuántos archivos \texttt{.v} tiene la verificación Coq de sotFS?
Liste sus nombres.
\end{ejercicio}

\begin{ejercicio}[\dificultad{$\star$}]
\label{ej:coq-2}
¿Cuáles son las cinco cláusulas de \texttt{WellFormed}? ¿Cuál falta
respecto del runtime?
\end{ejercicio}

\begin{ejercicio}[\dificultad{$\star\star$}]
\label{ej:coq-3}
Lea el \texttt{Theorem create\_preserves\_WellFormed} en
\citaCoq{DpoCreate.v} línea 335. Identifique los cinco lemmas
auxiliares que se invocan y la línea donde cada uno se prueba.
\end{ejercicio}

\begin{ejercicio}[\dificultad{$\star\star$}]
\label{ej:coq-4}
¿Por qué un \texttt{Admitted} no interrumpe la compilación de Coq?
Discuta los pros y contras de esa decisión de diseño.
\end{ejercicio}

\begin{ejercicio}[\dificultad{$\star\star$}]
\label{ej:coq-5}
Enuncie en Coq el invariante \texttt{NoHardLinkToDir} faltante. Pista:
una proposición \texttt{forall ce1 ce2 \textbackslash in g.edges,
ce1.tgt is a directory inode -> ce1.name = ".." \textbackslash/
ce2.name = "."}.
\end{ejercicio}

\begin{ejercicio}[\dificultad{$\star\star$}]
\label{ej:coq-6}
La trampa~\ref{tr:coq-admitted} dice que cerrar los tres
\texttt{Admitted} requiere ``re-probar'' los otros teoremas con la
invariante nueva. ¿Cuáles teoremas exactamente hay que tocar? Pista:
recorrer la tabla~\ref{tab:coq-files}.
\end{ejercicio}

\begin{ejercicio}[\dificultad{$\star\star\star$}]
\label{ej:coq-extension}
Formalice \texttt{NoHardLinkToDir} en Coq como una sexta cláusula de
\texttt{WellFormed} y pruebe que \opcreate{} la preserva. Para
\opmkdir, escriba el statement (no es necesario completar la prueba).
\end{ejercicio}

\begin{ejercicio}[\dificultad{$\star\star\star$}]
\label{ej:coq-extract}
Investigue cómo se haría la extracción Coq $\to$ Rust de las funciones
\texttt{create}, \texttt{mkdir}, \dots. ¿Qué plugins existen? ¿Cuál
sería el costo de mantenibilidad de tener dos implementaciones (Coq y
Rust) en sincronía vs.\ una sola (Rust) con prueba externa?
\end{ejercicio}
